\documentclass[12pt,a4paper]{article}
\usepackage[utf8]{inputenc}
\usepackage{graphicx}
\usepackage[none]{hyphenat}
\usepackage{hyperref}
\usepackage{xcolor}
\usepackage{pgfplots}
\usepackage{pgfplotstable}
\pgfplotsset{compat=1.17}

\definecolor{darkgreen}{RGB}{0, 150, 0}
\pgfplotstableread[col sep=comma]{../analysis_results.csv}\datatable

\begin{document}
	\begin{titlepage}
		
		\begin{center}
			\includegraphics[width=0.5\textwidth]{QUT.jpg}\\
			[0.03\textheight]  
			\Large\textbf{Bachelor of IT (Computer Science)}\\
			\Large\textbf{Assignment 2}\\
			\large\textbf{CAB301 - Algorithms and Complexity}\\
			[0.02\textheight]
			\large\textsl{Dane Madsen}\\
			\large\textsl{n10983864@qut.edu.au}
		\end{center}
		
	\end{titlepage}
	\tableofcontents
	\newpage
	
	\section{NoDVDs Method Design and Analysis}
		\subsection{Algorithm Design}
			This method calculates the total number of DVDs in a MovieCollection 
			through a in-order node tree traversal using a stack of BTreeNodes. Firstly, 
			it initializes the DVD count to be 0, creates an empty stack of BTreeNodes 
			and starts traversing from the root. Next, an outer while loop starts and 
			continues until the current BTreeNode is null and the stack count is equal 
			to 0. Inside the outer while loop, the inner while loop traverses the left 
			subtree, pushing left child BTreeNodes onto the stack. After traversing the 
			left subtree, the method pops the first movie from the top of the stack, adds 
			its total copies to the running total, and moves to the right child of the popped 
			movie. This process visits each node, cumulatively adding the total copies to 
			the count. Once the traversal finishes, it returns the total DVDs.\\
		
			\textbf{ALGORITHM} \textit{NoDVDs()}\\
			\null\hspace{1cm}// Calculates the total number of DVDs\\
			\null\hspace{1cm}// Returns the total number of DVDs in the MovieCollection\\
			\null\hspace{1cm}\textit{totalDVDs} $\leftarrow 0$\\
			\null\hspace{1cm}\textit{stack} $\leftarrow$ an empty BTreeNodes stack\\
			\null\hspace{1cm}\textit{curr} $\leftarrow$ root node of the MovieCollection\\
			\null\hspace{1cm}\textbf{while} \textit{curr} $\neq$ null \textbf{or} 
			\textit{stack} \textbf{is not} empty\\
			\null\hspace{2cm}\textbf{while} \textit{curr} $\neq$ null\\
			\null\hspace{3cm}\textit{stack.push(curr)}\\
			\null\hspace{3cm}\textit{curr} $\leftarrow$ \textit{curr.LChild}\\
			\null\hspace{2cm}\textit{curr} $\leftarrow$ \textit{stack.pop()}\\
			\null\hspace{2cm}\textit{totalDVDs} $\leftarrow$ \textit{totalDVDs} + 
			\textit{curr.Movie.TotalCopies}\\
			\null\hspace{2cm}\textit{curr} $\leftarrow$ \textit{curr.RChild}\\
			\null\hspace{1cm}\textbf{return} \textit{totalDVDs}

		\newpage

		\subsection{Algorithm Analysis}
			To perform an emperical analysis of the NoDVDs method, I created 20 
			MovieCollections of exponentially increasing sizes (from 1 to 524388 
			movies which is roughly close to the amount of movies on IMDb). Then, 
			for each of these MovieCollections, I ran the NoDVDs method 10 times 
			and averaged the time taken to run the method. The results of this 
			analysis are shown in Figure 1.\\

			\begin{figure}[h]
				\centering
				\begin{tikzpicture}
					\begin{axis}[
						title={Average Execution Time of NoDVDs()},
						xlabel={Collection Size},
						ylabel={Execution Time (ms)},
						xmode=log,
						ymode=log,
						log ticks with fixed point,
						grid=major,
						xmin=1, xmax=524288,
						ymin=0.0001, ymax=20,
						width=\textwidth,
						height=8cm
					]
					\addplot[mark=*,blue] table[x=MovieCollections,y=AveragedExecutionTime] from \datatable;
					\end{axis}
				\end{tikzpicture}
				\caption{Average Execution Time of NoDVDs()}
			\end{figure}
			
			As Figure 1 shows, though some of the first few MovieCollections have 
			timings that are outliers, after the third MovieCollection (which has 8 values), the average 
			execution times of the NoDVDs method increase linearly. As such the closest efficiency 
			classification for this method would be $O(n)$, where $n$ is the number of movies in the
			given MovieCollection. This hypothesis is can be supported by ploting the $T/n$ vs. $n$ as 
			shown in Figure 2.\\

			\newpage

			\begin{figure}[h]
				\centering
				\begin{tikzpicture}
				  \begin{axis}[
					title={Average Execution Time of NoDVDs() / n},
					xlabel={MovieCollection ID (n)},
					ylabel={Execution Time (T) / n [ms]},
					xmode=log,
					ymode=log,
					log ticks with fixed point,
					grid=major,
					width=\textwidth,
					height=8cm
				  ]
				  \addplot[mark=*,blue] table[x=MovieCollections,y=AveragedTN] from \datatable;
				  \end{axis}
				\end{tikzpicture}
				\caption{Plot of MovieCollection ID (n) vs. Execution Time (T) / n}
				\label{fig:execution_time_over_n_plot}
			\end{figure}
			
			Similar to Figure 1, Figure 2 shows that the first few MovieCollections have 
			$T/n$ values that are outliers, but again after the third MovieCollection, the $T/n$ 
			values become quite stable which supports the hypothesis that the algorithm used 
			for the NoDVDs method has an efficiency classification of $O(n)$.\\
		
		\newpage
	
	\section{Testing}
		\subsection{CompareTo Method Testing}
			To test the CompareTo method I created three XUnit tests. The first test 
			ensures that the method returns -1 when the movie CompareTo is being called 
			on is alphabetically less than the movie being provided to CompareTo as a 
			parameter. It achieves this by creating two movies, one with the letter A as 
			a title and the other with the letter B as a title. Then the unit test calls 
			CompareTo on the first movie (with the letter A) and passes the second movie 
			(with the letter B) as a parameter and asserts that the returned value is equal 
			to -1.\\
			\\
			The second test ensures that the CompareTo method returns 0 when the movie 
			CompareTo is being called on is equal to the movie being provided to CompareTo 
			as a parameter. It achieves this by creating two movies, both with the letter A 
			as a title. Then the unit test calls CompareTo on the first movie and passes the 
			second movie as a parameter, then it asserts that the returned value is equal to 
			0.\\
			\\
			The third and final test ensures that the CompareTo method returns 1 when the 
			movie CompareTo is being called on is alphabetically greater than the movie 
			being provided to CompareTo as a parameter. It achieves this by creating two 
			movies, one with the letter B as a title and the other with the letter A as a 
			title (the same as the first test only in reverse). Then the unit test calls 
			CompareTo on the first movie (with the letter B) and passes the second movie 
			(with the letter A) as a parameter and asserts that the returned value is equal 
			to 1.\\
			
		\newpage

		\subsection{ToString Method Testing}
			To test the ToString method I created three XUnit tests. The first test 
			ensures that the method returns the correct string when the movie has 
			only one copy. It achieves this by creating a movie with the parameters: 
			"A", Action, G, 1, 1. Then the unit test calls ToString on the movie and 
			asserts that the returned value is equal to the expected string. This test 
			can also be considered the minimum boundary test.\\
			\\
			The second test ensures that the ToString method returns the correct 
			string when the movie has more than one copy. It achieves this by creating 
			a movie with the parameters: "Sophies Choice", Drama, M, 151, 12137123. 
			Then the unit test calls ToString on the movie and asserts that the returned 
			value is equal to the expected string.\\
			\\
			The third and final test ensures that the ToString method returns the 
			correct string when the movie has a very large number of copies. It 
			achieves this by creating a movie with the parameters: "ZZZZZZZZZ", 
			Western, M15Plus, and the duration and totalcopies both 2147483647. 
			Then the unit test calls ToString on the movie and asserts that the 
			returned value is equal to the expected string. This test can also 
			be considered the maximum boundary test.\\

		\subsection{IsEmpty Method Testing}
			To test the IsEmpty method I created two XUnit tests. The first test 
			ensures that the method returns true when the MovieCollection is empty. 
			It achieves this by creating an empty MovieCollection and asserting that 
			when the IsEmpty method is called on the MovieCollection, it returns 
			true.\\
			\\
			The second test ensures that the IsEmpty method returns false when the 
			MovieCollection is not empty. It achieves this by again creating a 
			MovieCollection, but this time populating it with a movie. Then	the 
			unit test asserts that when the IsEmpty method is called on the 
			MovieCollection, it returns false.\\
		
		\newpage

 		\subsection{Insert Method Testing}
			To test the Insert method I created four XUnit tests. The first test 
			ensures that the method returns true when the movie being inserted has 
			minimum boundary values. It achieves this by creating a MovieCollection 
			and a movie with the parameters: "A", Action, G, 1, 1. Then the unit 
			test calls Insert on the MovieCollection and passes the movie as a 
			parameter and asserts that the returned value is equal to true.\\
			\\
			The second test ensures that the method returns true when the movie 
			being inserted does not already exist in the MovieCollection. It 
			achieves this by creating a MovieCollection and a movie with the 
			parameters: "Sophies Choice", Drama, M, 151, 12137123. Then the unit 
			test calls Insert on the MovieCollection and passes the movie as a 
			parameter and asserts that the returned value is equal to true.\\
			\\
			The third test ensures that the Insert method fuctions correctly when 
			the movie being inserted has maximum boundary values. It achieves this 
			by creating a MovieCollection and a movie with the parameters: "ZZZZZZZZZ", 
			Western, M15Plus, and the duration and totalcopies both 2147483647. Then 
			the unit test calls Insert on the MovieCollection and passes the movie 
			as a parameter and asserts that the returned value is equal to true.\\
			\\
			The fourth test ensures that the Insert method returns false when the 
			movie being inserted already exists in the MovieCollection. It achieves 
			this by creating a MovieCollection and a movie with the parameters: 
			"A", Action, G, 1, 1. Then the unit test inserts the movie into the 
			MovieCollection once, ensures that the returned value is true, and 
			then inserts the movie into the MovieCollection again and asserts 
			that the returned value is false.\\

		\subsection{Delete Method Testing}
			To test the Delete method I created three XUnit tests. The first test 
			ensures that the method returns true when the movie has successfully 
			been deleted. It achieves this by creating a MovieCollection and a 
			movie with the parameters: "A", Action, G, 1, 1. Then the unit test 
			inserts the movie into the MovieCollection, ensuring that the returned 
			value is true, and then, deletes the same movie from the MovieCollection 
			and asserts that the returned value is again true.\\
			\\
			\\
			The second test ensures that the Delete method returns false when the 
			MovieCollection is empty. It achieves this by creating an empty 
			MovieCollection and a movie with the parameters: "A", Action, G, 1, 1. 
			Then the unit test calls Delete on the MovieCollection and passes the 
			movie as a parameter and asserts that the returned value is equal to 
			false.\\
			\\
			The third and final test ensures that the Delete method returns false 
			when the movie being deleted does not exist in the MovieCollection. 
			It achieves this by creating a MovieCollection and a movie with the 
			parameters: "A", Action, G, 1, 1. Then the unit test inserts the movie 
			into the MovieCollection, ensuring that the returned value is true, 
			and then, attempts to delete a different movie from the MovieCollection 
			(with the parameters: "B", Action, G, 1, 1) and asserts that the returned 
			value is false.\\

		\subsection{Search Method Testing}
			To test the Search method I created two XUnit tests. The first test 
			ensures that the method returns the correct movie when the movie being 
			searched for exists in the MovieCollection. It achieves this by first 
			creating a MovieCollection and then a movie with the parameters: "A", 
			Action, G, 1, 1. Then the unit test inserts the movie into the 
			MovieCollection. Finally, the unit test calls Search on the 
			MovieCollection and passes "A" as a parameter (the title of the movie) 
			and asserts that the value returned is equal to the movie that was 
			inserted.\\
			\\
			The second test ensures that the Search method returns null when the 
			movie being searched for does not exist in the MovieCollection. It 
			achieves this by creating a MovieCollection and then immediately 
			calling Search on the MovieCollection with "A" as a parameter. Then 
			the unit test asserts that the value returned is equal to null.\\

		\subsection{NoDVDs Method Testing}
			To test the NoDVDs method I created three XUnit tests. The first test 
			ensures that the method returns 0 when the MovieCollection is empty. 
			It achieves this by creating an empty MovieCollection and asserting that 
			when the NoDVDs method is called on the MovieCollection, it returns 0. 
			This can also be considered the minimum boundary test.\\ 
			\\
			The second test ensures that NoDVDs returns the correct number of DVDs under 
			normal circumstances. To achieve this the unit test creates a MovieCollection 
			and then populates it with two movies each with 1000 totalDVDs. Then, the 
			unit test asserts that the the returned value of NoDVDs is equal to 2000.\\
			\\
			The third and final test ensures that the NoDVDs method returns the correct 
			number of DVDs when the amount of DVDs is very large. To achieve this, the 
			unit test creates a MovieCollection and then populates it with two movies 
			each with half the maximum integer value (2147483647) minus one (1073741823, 
			because the maxint is odd) as their totalDVDs. Then, the unit test asserts 
			that the returned value of NoDVDs is equal to the maximum integer value 
			minus one.\\

		\subsection{ToArray Method Testing}
			To test the ToArray method I created two XUnit tests. The first test ensures 
			that the ToArray method returns an empty array when the MovieCollection is empty. 
			It achieves this by creating an empty MovieCollection and then asserting that the 
			returned array is equal to a user defined empty array.\\
			\\
			The second test ensures that the method returns the correct array when the 
			ToArray method is called. It achieves this by creating a MovieCollection 
			and then populating it with two movies. Then, the unit test asserts that 
			the returned array is equal to a user defined array with the same movies.\\

		\subsection{Clear Method Testing}
			To test the Clear method I created two XUnit tests. The first test ensures 
			that the method functions when the MovieCollection is empty. It achieves 
			this by creating an empty MovieCollection, asserting that it is empty, and 
			then calling Clear on the MovieCollection and asserting again that it is 
			empty.\\
			\\
			The second test ensures that the method functions when the MovieCollection 
			is not empty. It achieves this by creating a MovieCollection and then 
			populating it with a movie. Then, the unit test asserts that the MovieCollection 
			is not empty, calls Clear on the MovieCollection, and asserts that the 
			MovieCollection is empty.\\

\end{document}